\begin{abstract}
Following public emergencies, social media becomes a core arena for both information dissemination and emotional expression. Understanding the dynamics of public opinion is therefore crucial for crisis management. Existing large language model (LLM)-based multi-agent simulation frameworks mainly focus on rumor propagation and adversarial scenarios, while still lacking systematic modeling of official information intervention and cognitive reflection under real crisis conditions. This paper develops a time-aware LLM-driven multi-agent public opinion simulation system for public emergencies, integrating an official publisher mechanism and a cognitive reflection module. We validate the framework on three real cases---the 2022 China Eastern MU5735 crash, the 2023 Guangzhou BMW crash incident, and the 2023 Cathay Pacific discrimination incident---through four ablation settings and four quantitative metrics (\(\Delta\)Bias, \(\Delta\)Div, Corr, and DTW\_norm). By conducting multi-case empirical validation, this study extends LLM-based public opinion simulation from single-case demonstration to cross-event generalization analysis, and provides a validated simulation platform for counterfactual policy experiments.
This abstract is a placeholder. Revise it after the remaining sections are drafted and the target format is confirmed.
\end{abstract}
